% ============================================================
% LAMPIRAN B - KATALOG PERINTAH VERIFIKASI USE CASE
% Katalog perintah yang dapat dijalankan ulang (read-only terhadap
% kluster langsung) untuk membuktikan tiap skenario SK-F/SK-N pada
% Bab VI. Dikelompokkan mengikuti tujuan platform T-1..T-4 agar rantai
% T-N <-> SK <-> KF/KNF tetap satu lawan satu. Bersumber dari direktori
% experiments/ (sejajar platform/ dan use-case-crypto/); rincian argumen
% tiap skenario ada pada README.md di setiap subdirektori.
% ============================================================
\cleardoublepage
\chapter{KATALOG PERINTAH VERIFIKASI USE CASE}
\label{appx:verifikasi}

Lampiran ini memuat katalog perintah yang dapat dijalankan ulang untuk membuktikan jalur ujung ke ujung pada \textit{use case} pasar aset kripto. Setiap perintah dijalankan secara manual dan bersifat \textit{read-only} terhadap kluster yang sedang berjalan, sehingga tidak menjadi bagian dari \texttt{make phase-full} maupun rekonsiliasi Argo CD dan tidak mengubah sistem yang diuji, sejalan dengan metode evaluasi pada Bab~\ref{chap:evaluasi}. Berkas skenario berada pada direktori \texttt{experiments/} yang sejajar dengan \texttt{platform/} dan \texttt{use-case-crypto/}; rincian argumen tiap skenario tersedia pada berkas \texttt{README.md} di masing-masing subdirektori. Prasyarat eksekusi adalah platform telah terpasang (\texttt{make phase-full}), \textit{use case} telah ter-\textit{deploy} (\texttt{make usecase-crypto-build \&\& make usecase-crypto-up}), dan seluruh \textit{pod} berstatus \textit{Running}. Sebelum menjalankan skenario apa pun, konfigurasi domain dimuat dan \textit{namespace} pembangkit beban dibuat melalui preambel berikut.

\begin{lstlisting}[caption={Preambel: pemuatan konfigurasi domain dan \textit{namespace} pembangkit beban},basicstyle=\ttfamily\footnotesize,breaklines=true,captionpos=t]
# Muat blok DOMAIN dari satu titik konfigurasi (entitas, tabel medali,
# endpoint Feast/Qdrant/ClickHouse/DataHub, knob beban). Ganti use case =
# ganti berkas ini, bukan kode skenario.
set -a; . "${EXPERIMENTS_CONFIG:-experiments/config.crypto.env}"; set +a

# Namespace terpisah untuk pembangkit beban (k6, chaos) — bukan namespace produksi.
kubectl apply -f experiments/namespace.yaml
\end{lstlisting}

\section{Tujuan T-1: Arsitektur Terintegrasi dan GitOps}
\label{sec:appx-t1}

Kelompok T-1 membuktikan pemasangan deklaratif dari satu sumber Git, otomasi siklus hidup model setara MLOps Level 2, serta atribut kualitas lintas komponen. Perintahnya dirangkum pada Listing~\ref{lst:cat-t1}.

\begin{lstlisting}[caption={Perintah verifikasi tujuan T-1},label={lst:cat-t1},basicstyle=\ttfamily\footnotesize,breaklines=true,captionpos=t]
# SK-F-01 (KF-01) — GitOps reconciliation: ubah definisi feature view, dorong ke
# Gitea, amati Argo CD menyinkronkan, lalu daftarkan feature view aktif.
$EDITOR use-case-crypto/scripts/libs/feature_store/definitions.py
git -C use-case-crypto commit -am "feat: feature view" && git -C use-case-crypto push gitea
argocd app sync crypto-use-case-local && uv run feast feature-views list

# SK-F-12 (KF-12) — konfigurasi deklaratif: ubah jumlah replica di overlay, dorong,
# sinkronkan, lalu pastikan tidak ada drift.
argocd app diff crypto-use-case-local
kubectl -n use-case-crypto get deploy

# SK-F-06 (KF-06) — promosi model kanari + rollback model bermasalah (Argo Rollouts).
kubectl -n use-case-crypto get rollouts
kubectl argo rollouts -n use-case-crypto get rollout <rollout> --watch   # promote | abort

# SK-F-13 (KF-13) — antarmuka terpadu: satu sesi UV memanggil SDK Feast, MLflow,
# dan klien KServe pada satu notebook (rincian: experiments/feast/README.md).
uv run feast feature-views list                          # Feast SDK
kubectl get inferenceservice -A | grep predictor         # KServe (crypto-predictor)
# MLflow tracking client dipanggil dari notebook yang sama (konsol via ingress).

# SK-F-14 (KF-14) — manajemen sumber daya Kueue: jenuhkan ClusterQueue; job di luar
# anggaran berstatus Pending.
kubectl get clusterqueue,localqueue,workloads -A

# SK-N-03 (KNF-03) — skalabilitas horizontal HPA + KEDA pada beban naik.
k6 run experiments/load/hpa-keda-rampup.js
kubectl -n use-case-crypto get hpa,scaledobject,pods -w

# SK-N-04 (KNF-04) — keandalan: injeksi fault pod + jaringan (RTO<5m, RPO<1m).
envsubst < experiments/chaos/pod-chaos.yaml | kubectl apply -f -
kubectl apply -f experiments/chaos/network-chaos.yaml

# SK-N-08 (KNF-08) — ekstensibilitas: ganti runtime serving / online store lewat
# CRD + overlay Kustomize (mis. Seldon sebagai alternatif), perubahan terisolasi.
kubectl apply -k use-case-crypto/manifests/overlays/local

# SK-N-09 (KNF-09) — kemudahan penggunaan: akses seluruh konsol lewat satu sesi
# identitas dex + oauth2-proxy (SSO, tanpa login ulang).
kubectl get pods -A | grep -E 'dex|oauth2-proxy'

# SK-N-10 (KNF-10) — efisiensi sumber daya: rasio kompresi + atribusi biaya OpenCost.
kubectl get pods -A | grep -i opencost

# SK-N-11 (KNF-11) — reproduksibilitas: terapkan ulang seluruh manifes pada node
# k3s baru dari satu sumber Git.
make phase-full

# SK-N-12 (KNF-12) — cakupan uji (>80%) + waktu konvergensi kluster bersih.
make usecase-crypto-test
time make phase-full
\end{lstlisting}

\section{Tujuan T-2: Tata Kelola Data}
\label{sec:appx-t2}

Kelompok T-2 membuktikan katalog metadata, \textit{lineage}, kontrak kualitas, observabilitas, dan keamanan berjalan otomatis sepanjang \textit{pipeline}. Perintahnya dirangkum pada Listing~\ref{lst:cat-t2}.

\begin{lstlisting}[caption={Perintah verifikasi tujuan T-2},label={lst:cat-t2},basicstyle=\ttfamily\footnotesize,breaklines=true,captionpos=t]
# SK-F-08 (KF-08) — lineage end-to-end di DataHub (topik sumber -> prediksi).
experiments/lineage/lineage-walk.sh discover
experiments/lineage/lineage-walk.sh walk '<urn>'         # urn dari hasil discover

# SK-F-09 (KF-09) — pelacakan eksperimen: dua run pelatihan + replay dari snapshot
# (delta metrik <=1%); diamati pada MLflow.
kubectl get pods -A | grep -i mlflow                     # konsol MLflow via ingress

# SK-F-10 (KF-10) — registri + audit trail: transisi status model
# none -> staging -> production -> archived (MLflow Model Registry).
kubectl get pods -A | grep -i mlflow

# SK-N-05 (KNF-05) — kontrak kualitas: checkpoint Great Expectations + injeksi
# pelanggaran skema/rentang/kelengkapan.
kubectl -n use-case-crypto get cronjob | grep -iE 'quality|expect|ge-'
kubectl -n use-case-crypto create job ge-check --from=cronjob/<ge-checkpoint>

# SK-N-06 (KNF-06) — observabilitas: telusuri satu trace OpenTelemetry di Tempo,
# ingestasi -> prediksi pada satu trace ID (kueri lewat Grafana).
kubectl get pods -A | grep -E 'tempo|grafana'

# SK-N-07 (KNF-07) — keamanan: pindai cluster dengan Trivy + periksa TLS 1.3/AES-256.
kubectl get vulnerabilityreports -A
testssl.sh <public-endpoint>
\end{lstlisting}

\section{Tujuan T-3: Deteksi \textit{Drift} dan Pelatihan Ulang}
\label{sec:appx-t3}

Kelompok T-3 membuktikan deteksi \textit{drift} terotomasi memicu pelatihan ulang dan penyusunan fitur pelatihan menjaga \textit{point-in-time correctness}. Perintahnya dirangkum pada Listing~\ref{lst:cat-t3}.

\begin{lstlisting}[caption={Perintah verifikasi tujuan T-3},label={lst:cat-t3},basicstyle=\ttfamily\footnotesize,breaklines=true,captionpos=t]
# SK-F-03 (KF-03) — point-in-time correctness: audit get_historical_features pada
# gold.fct_training_data dengan probe event_timestamp masa depan; setiap baris
# bermuatan informasi setelah timestamp acuan ditolak (zero leakage).
# Prosedur SDK Feast: experiments/feast/README.md. Cek substrat dulu:
experiments/etl/data-landed-check.sh

# SK-F-05 (KF-05) — otomasi pelatihan di Kubeflow Pipelines (baca Feast -> tulis MLflow).
uv run --with 'kfp[kubernetes]==2.16.0' use-case-crypto/pipelines/retraining_pipeline.py

# SK-F-07 (KF-07) — pemantauan + retrain otomatis: geser distribusi data nyata
# beberapa sigma, amati detektor drift + CronWorkflow retrain-on-drift.
uv run experiments/drift/inject_drift.py --sigma 2.0
kubectl -n use-case-crypto get cronworkflow,workflow
\end{lstlisting}

\section{Tujuan T-4: Layanan Fitur \textit{Dual-Store}}
\label{sec:appx-t4}

Kelompok T-4 membuktikan layanan fitur \textit{dual-store} menyajikan fitur tabular, agregat, dan vektor dengan kontrak konsisten antara pelatihan dan penyajian, beserta latensi dan \textit{throughput} jalur penyajian. Perintahnya dirangkum pada Listing~\ref{lst:cat-t4}.

\begin{lstlisting}[caption={Perintah verifikasi tujuan T-4},label={lst:cat-t4},basicstyle=\ttfamily\footnotesize,breaklines=true,captionpos=t]
# SK-F-02 (KF-02) — penyajian ganda: online non-null (Feast) vs offline (ClickHouse).
kubectl -n model-lifecycle port-forward svc/feast 6566:6566 &
FEAST_URL=http://localhost:6566 uv run experiments/feast/online-serving-check.py
experiments/etl/data-landed-check.sh

# SK-F-04 (KF-04) — pencarian vektor (Qdrant top-k recall + latensi).
k6 run experiments/load/vector-search-latency.js

# SK-F-11 (KF-11) — paritas batch (dbt volume_sma_20) vs stream (Flink secondary_avg)
# pada rata-rata volume 20-event yang sama; cek kesehatan substrat dulu.
experiments/etl/data-landed-check.sh
experiments/parity/parity-check.sh

# SK-N-01 (KNF-01) — latensi feature serving + vector search (p99 sub-detik).
k6 run experiments/load/feast-online-latency.js

# SK-N-02 (KNF-02) — throughput serving + stream; pantau consumer lag.
k6 run experiments/load/feast-throughput.js
# Setengah stream: kredensial SASL Kafka diambil langsung dari kluster (tidak
# dikomit). Topik crypto.validated; pantau metrik kafka_consumergroup_lag.
kafka-producer-perf-test --topic crypto.validated --num-records 1000000 \
  --record-size 256 --throughput -1 \
  --producer-props bootstrap.servers="$KAFKA_BOOTSTRAP" \
  --producer.config <sasl.properties>
\end{lstlisting}

Setiap perintah pada katalog ini berpasangan satu lawan satu dengan baris Tabel~\ref{tab:uji_penerimaan} dan Tabel~\ref{tab:evaluasi_target} pada Bab~\ref{chap:evaluasi}, sehingga klaim pemenuhan kebutuhan dapat ditelusuri langsung ke perintah yang menghasilkan buktinya. Untuk \textit{use case} lain, hanya berkas \texttt{config.crypto.env} yang disalin dan disesuaikan blok DOMAIN-nya; kode skenario tetap \textit{domain-agnostic} dan tidak perlu disentuh.
